% The manual's preamble: fonts, colours, page furniture, tables, code, figures, index, links.
% Built with LuaLaTeX (latexmk -lualatex); no shell escape is needed.

\usepackage{fontspec}
\setmainfont{LibertinusSerif}[
  Extension=.otf, UprightFont=*-Regular, BoldFont=*-Semibold,
  ItalicFont=*-Italic, BoldItalicFont=*-SemiboldItalic, Ligatures=TeX]
\setsansfont{LibertinusSans}[
  Extension=.otf, UprightFont=*-Regular, BoldFont=*-Bold, ItalicFont=*-Italic,
  Ligatures=TeX]
\setmonofont{Inconsolatazi4}[
  Extension=.otf, UprightFont=*-Regular, BoldFont=*-Bold, Scale=MatchLowercase]

\usepackage[english]{babel}
\usepackage{microtype}
\usepackage[dvipsnames,table]{xcolor}
\usepackage{graphicx}
\usepackage{amsmath}
\usepackage{array,booktabs,longtable,tabularx,multirow}
\usepackage{enumitem}
\usepackage{caption}
\usepackage{float}
\graphicspath{{screenshots/}}
% \screenshot[includegraphics options]{file}{caption}{label}: a screenshot of the running application
\newcommand{\screenshot}[4][]{%
  \begin{figure}[htbp]\centering
  \setlength{\fboxsep}{0pt}\setlength{\fboxrule}{.4pt}%
  \fcolorbox{rule}{white}{\includegraphics[width=\dimexpr\linewidth-1pt\relax,#1]{#2}}%
  \caption{#3}\label{#4}\end{figure}}
\usepackage{csquotes}
\usepackage[obeyspaces,spaces]{url}
\usepackage{listings}
\usepackage{tikz}
\usepackage{pgfplots}
\usepackage{pgfplotstable}
\usepackage[noautomatic]{imakeidx}
\usepackage{scrlayer-scrpage}

% ---------------------------------------------------------------------------- colours
% One palette for the whole book: cobalt for structure and links, signal orange for the one
% thing a figure is about, ink and slate for text, rule and paper for furniture.
\definecolor{cobalt}{HTML}{1F4FCC}
\definecolor{cobaltlight}{HTML}{E8EEFB}
\definecolor{sky}{HTML}{8FB3FF}
\definecolor{signal}{HTML}{E8641B}
\definecolor{signallight}{HTML}{FCEBDD}
\definecolor{ink}{HTML}{1C2230}
\definecolor{slate}{HTML}{5A6475}
\definecolor{rule}{HTML}{C9D0DC}
\definecolor{paper}{HTML}{F5F7FA}
\definecolor{moss}{HTML}{2E8B57}

% ---------------------------------------------------------------------------- page and headings
\setlength{\emergencystretch}{2em}
\clubpenalty=10000
\widowpenalty=10000
\displaywidowpenalty=10000
\addtokomafont{disposition}{\color{ink}}
\setkomafont{chapterprefix}{\normalfont\sffamily\Large\color{slate}}
\addtokomafont{chapter}{\color{cobalt}}
\setkomafont{part}{\normalfont\bfseries\fontsize{30}{36}\selectfont\color{cobalt}}
\setkomafont{partnumber}{\normalfont\sffamily\LARGE\color{slate}}
\setkomafont{descriptionlabel}{\normalfont\bfseries}
\addtokomafont{caption}{\small}
\addtokomafont{captionlabel}{\sffamily\bfseries\color{cobalt}}
\setkomafont{pageheadfoot}{\small\sffamily\color{slate}}
\setkomafont{pagenumber}{\small\sffamily\color{ink}}
\renewcommand*{\sectionformat}{\textcolor{slate}{\thesection\autodot}\enskip}
\renewcommand*{\subsectionformat}{\textcolor{slate}{\thesubsection\autodot}\enskip}
\RedeclareSectionCommand[beforeskip=-2.4ex plus -1ex minus -.2ex,afterskip=1.1ex plus .2ex]{section}
\RedeclareSectionCommand[beforeskip=-1.8ex plus -1ex minus -.2ex,afterskip=.7ex plus .2ex]{subsection}
% a signal rule closes every chapter heading
\renewcommand*{\chapterheadendvskip}{\par\nobreak\vspace{1.6ex}{\color{signal}\hrule height .9pt}\nobreak\vspace{4.5ex}}
\renewcommand*{\partpagestyle}{empty}
% \partintro{text}: the short preamble under a part title, after a signal rule
\newcommand{\partintro}[1]{\setpartpreamble[c][.78\textwidth]{\vspace*{3ex}{\centering\color{signal}\rule{3cm}{1.2pt}\par}\vspace{3.5ex}\small\noindent #1}}
\setkomafont{partentry}{\sffamily\bfseries\color{cobalt}}
\setkomafont{chapterentry}{\bfseries}
\setkomafont{chapterentrypagenumber}{\normalfont}

% running heads: the chapter on the inside, the section on the outside, a hairline under both
\pagestyle{scrheadings}
\clearpairofpagestyles
\ihead{\leftmark}
\ohead{\rightmark}
\ofoot*{\pagemark}
\automark[section]{chapter}
\KOMAoptions{autooneside=false}
\renewcommand*{\chaptermarkformat}{\thechapter\autodot\enskip}
\renewcommand*{\sectionmarkformat}{\thesection\autodot\enskip}
\KOMAoptions{headsepline=.4pt}
\setkomafont{headsepline}{\color{rule}}
\makeatletter\renewcommand\@pnumwidth{2.2em}\makeatother
\setcounter{tocdepth}{1}
\setcounter{secnumdepth}{2}
\setlist{itemsep=.2ex,topsep=.6ex}
\setlist[description]{style=nextline, font=\normalfont\bfseries}
\captionsetup{font=small,labelfont={sf,bf,color=cobalt},format=plain,justification=raggedright,singlelinecheck=false,skip=6pt}

\renewcommand{\topfraction}{.9}
\renewcommand{\bottomfraction}{.7}
\renewcommand{\textfraction}{.08}
\renewcommand{\floatpagefraction}{.75}
\setcounter{topnumber}{3}
\setcounter{bottomnumber}{2}
\setcounter{totalnumber}{5}
% floats on a float page sit at the top, not in the middle
\makeatletter\setlength{\@fptop}{0pt}\setlength{\@fpbot}{0pt plus 1fil}\makeatother

% ---------------------------------------------------------------------------- inline code and terms
\urlstyle{tt}
\makeatletter
\g@addto@macro{\UrlBreaks}{\do\_\do\-\do\.\do\=\do\,\do\:}
\makeatother
\Urlmuskip=0mu plus 1mu
% \code{...}: identifiers, flags, variables, routes and paths, set in monospace and broken
% after punctuation; for text with braces, %, # or \ use \texttt with escapes instead.
\DeclareUrlCommand\code{\urlstyle{tt}}
% \term{word}: a defined term, set in italics and put in the index
\newcommand{\term}[1]{\emph{#1}\index{#1}}
% the three kinds of number (About this manual, "How numbers are labelled")
\newcommand{\measured}{\textsf{\small measured}}
\newcommand{\arith}{\textsf{\small arithmetic}}
% \hd{...}: a table header cell
\newcommand{\hd}[1]{{\sffamily\color{slate}#1}}

% ---------------------------------------------------------------------------- tables
\renewcommand{\arraystretch}{1.15}
\arrayrulecolor{rule}
\newcolumntype{L}[1]{>{\raggedright\arraybackslash}p{#1}}
\newcolumntype{Y}{>{\raggedright\arraybackslash}X}
% reftable: a long reference table with three columns, a caption, a label and a header
\newenvironment{reftable}[3]{%
  \footnotesize
  \setlength{\LTpre}{.8ex}\setlength{\LTpost}{.8ex}%
  \longtable{@{}L{.29\linewidth}L{.19\linewidth}L{.445\linewidth}@{}}
  \caption{#1}\label{#2}\\
  \toprule #3\\ \midrule \endfirsthead
  \multicolumn{3}{@{}l}{\textit{\small continued}}\\ \toprule #3\\ \midrule \endhead
  \bottomrule \endlastfoot
}{\endlongtable}
% httptable: the HTTP routes, three columns: method and path, who may call, what it does
\newenvironment{httptable}[2]{%
  \footnotesize
  \setlength{\LTpre}{.8ex}\setlength{\LTpost}{.8ex}%
  \longtable{@{}L{.31\linewidth}L{.09\linewidth}L{.55\linewidth}@{}}
  \caption{#1}\label{#2}\\
  \toprule \hd{route} & \hd{who} & \hd{what, body, answer}\\ \midrule \endfirsthead
  \multicolumn{3}{@{}l}{\textit{\small continued}}\\ \toprule \hd{route} & \hd{who} & \hd{what, body, answer}\\ \midrule \endhead
  \bottomrule \endlastfoot
}{\endlongtable}
% reftabletwo: the same with two columns, a name and what it is
\newenvironment{reftabletwo}[3]{%
  \footnotesize
  \setlength{\LTpre}{.8ex}\setlength{\LTpost}{.8ex}%
  \longtable{@{}L{.33\linewidth}L{.62\linewidth}@{}}
  \caption{#1}\label{#2}\\
  \toprule #3\\ \midrule \endfirsthead
  \multicolumn{2}{@{}l}{\textit{\small continued}}\\ \toprule #3\\ \midrule \endhead
  \bottomrule \endlastfoot
}{\endlongtable}

% ---------------------------------------------------------------------------- code listings
% One style for every block of code, commands, output and status lines.
\lstdefinestyle{tandem}{
  basicstyle=\ttfamily\small,
  columns=fullflexible, keepspaces=true,
  breaklines=true, breakatwhitespace=false,
  postbreak=\mbox{\textcolor{slate}{$\hookrightarrow$}\space},
  frame=leftline, framerule=1.2pt, rulecolor=\color{cobalt},
  backgroundcolor=\color{paper},
  xleftmargin=1.2em, framexleftmargin=.6em, aboveskip=1.2ex, belowskip=1.2ex,
  showstringspaces=false, upquote=true, extendedchars=true,
  literate={›}{{\guilsinglright}}1 {µ}{{\textmu}}1 {~}{{\textasciitilde}}1,
  commentstyle=\color{slate},
  keywordstyle=\color{cobalt},
  stringstyle=\color{ink},
}
\lstset{style=tandem}
\lstdefinelanguage{shellx}{morekeywords={export,set,bash,cd,pip,curl,hf,python,python3,docker,sudo,tandem-dash,git,vllm,flock,printf,echo,grep,security,rm,cp,cat,sqlite3,npm,npx},
  sensitive=true, morecomment=[l]{\#}, morestring=[b]"}
\lstnewenvironment{shell}{\lstset{language=shellx}}{}
\lstnewenvironment{plain}{\lstset{language={}}}{}
\lstnewenvironment{lines}{\lstset{language={}}}{}
\lstnewenvironment{jsonc}{\lstset{language={},morestring=[b]"}}{}

% ---------------------------------------------------------------------------- figures
\usetikzlibrary{positioning,arrows.meta,shapes.geometric,shapes.misc,fit,backgrounds,calc,
  decorations.pathreplacing,patterns,matrix,chains}
\pgfplotsset{compat=1.18}
\tikzset{
  >={Stealth[length=2.2mm]},
  box/.style={draw=cobalt, fill=cobaltlight, rounded corners=2pt, align=center,
    font=\sffamily\small, inner sep=4pt, minimum height=8mm},
  hot/.style={draw=signal, fill=signallight, rounded corners=2pt, align=center,
    font=\sffamily\small, inner sep=4pt, minimum height=8mm},
  quiet/.style={draw=rule, fill=paper, rounded corners=2pt, align=center,
    font=\sffamily\small, inner sep=4pt, minimum height=8mm},
  lbl/.style={font=\sffamily\scriptsize, text=slate, align=center},
  flow/.style={->, thick, draw=ink},
  dflow/.style={->, thick, draw=slate, dashed},
  group/.style={draw=rule, rounded corners=4pt, inner sep=6pt},
  % the instrument-readout pieces of the dashboard figures
  panel/.style={draw=rule, fill=white, rounded corners=3pt, inner sep=5pt},
  readout/.style={font=\sffamily\bfseries\Large, text=ink},
  unit/.style={font=\sffamily\scriptsize, text=slate},
}
% ---------------------------------------------------------------------------- the mark
% The TandemLLM mark (Scope: a dial, the StairCut step trace inside it, the orange now-rate
% breaking out on the right), drawn in TikZ from the geometry of docs/media/logo/mark.svg so the
% build needs no SVG converter. Its own colours, not the book's: brandcobalt and brandsignal.
\definecolor{brandcobalt}{HTML}{3D63FF}
\definecolor{brandsignal}{HTML}{FF6B1F}
% \tandemmark[<unit>]{<dial and trace>}{<now-rate>}: the 64-unit grid at one unit = <unit>;
% the box is the mark's ink, units 2 to 62.
\newcommand{\tandemmark}[3][.25mm]{%
  \begin{tikzpicture}[x=#1, y=-#1]
    \useasboundingbox (2,2) rectangle (62,62);
    \draw[draw=#2, line width=\dimexpr#1*4\relax] (32,32) circle[radius=26];
    \draw[draw=#2, line width=\dimexpr#1*6\relax, line cap=butt, line join=round,
      rounded corners=\dimexpr#1*3/2\relax]
      (8,42) -- (22,42) -- (22,30) -- (38,30) -- (38,18) -- (41,18);
    \draw[draw=#3, line width=\dimexpr#1*6\relax, line cap=round] (40,18) -- (60,18);
  \end{tikzpicture}}
% built once, outside any other picture (a tikzpicture inside the cover's overlay would inherit it)
\newsavebox{\tandemcoverbox}
\newsavebox{\tandempartbox}
\AtBeginDocument{%
  \sbox{\tandemcoverbox}{\tandemmark[.32mm]{brandcobalt}{brandsignal}}%
  \sbox{\tandempartbox}{\tandemmark[.14mm]{slate}{slate}}}
% a small mono mark above "Part I" on each part page
\renewcommand*{\partformat}{\usebox{\tandempartbox}\\[2.2ex]\partname~\thepart\autodot}

\pgfplotsset{
  tandem/.style={
    width=\linewidth, height=6.2cm,
    axis lines*=left, grid=major, grid style={draw=rule!60},
    tick label style={font=\sffamily\footnotesize}, label style={font=\sffamily\small},
    legend style={font=\sffamily\footnotesize, draw=none, fill=none},
    every axis plot/.append style={thick},
  },
  % a single stacked bar: no y axis, no empty-range warning
  onebar/.style={ytick=\empty, axis y line=none, ymin=-1, ymax=1, enlarge y limits=false},
}

% ---------------------------------------------------------------------------- index and links
\makeindex[intoc, columns=2, title=Index, options=-s tandem.ist]
\usepackage[colorlinks=true,linkcolor=cobalt,urlcolor=cobalt,citecolor=cobalt,
  bookmarksnumbered=true,bookmarksopen=true,bookmarksopenlevel=0,unicode=true,
  pdftitle={TandemLLM Manual},pdfauthor={The TandemLLM authors},
  pdfsubject={An inference engine for Qwen3.8-27B on one DGX Spark, its conductor and its dashboard},
  pdfkeywords={speculative decoding, inference, NVFP4, DGX Spark, StairCut}]{hyperref}
\usepackage{bookmark}
\newcommand{\version}{0.13.0}
\newcommand{\manualdate}{October 2026}

% "see" entries of the index
\index{lookup store|see{suffix store}}
\index{n-gram drafter|see{lookup drafter}}
\index{KV cache|see{resident prefix, state store}}
\index{queue|see{admission}}
